\begin{textsample}{POS Dim 6 – to – Score 9.00 – to/to\_em\_14.txt}  \label{ex:f6_pos_174}
APRESENTAÇÃO O Caderno 3, denominado Itinerários \textbf{Formativos} - \textbf{Trilhas} de Aprofundamento, busca apoiar o trabalho dos docentes, contribuir com a organização pedagógica das unidades escolares e fomentar a diversificação curricular, com vistas à formação integral dos estudantes.
A arquitetura curricular do Ensino Médio do território do Tocantins está organizada em Formação Geral Básica e Itinerários \textbf{Formativos}.
Os Itinerários \textbf{Formativos} são compostos por : \textbf{Trilhas} de Aprofundamentos e/ou de Formação Técnica Profissional, Eletivas e Projeto de Vida.
Este é o Caderno 3 – que aborda os Itinerários \textbf{Formativos} – \textbf{Trilhas} de Aprofundamento e está subdividido em quatro blocos : 3.1 - \textbf{Trilhas} de Aprofundamento de Linguagens e suas Tecnologias ; 3.2 - \textbf{Trilhas} de Aprofundamento de Matemática e suas Tecnologias ; 3.3 - \textbf{Trilhas} de Aprofundamento de Ciências Humanas e Sociais Aplicadas ; 3.4 - \textbf{Trilhas} de Aprofundamento de Ciências da Natureza e suas Tecnologias.
No intuito de atender as atualizações dos marcos legais do Ensino Médio que promoveram mudanças, a organização curricular passa a ser composta por duas partes indissociáveis : a Formação Geral Básica ( FGB ), com até 1.800 horas de trabalho pedagógico, e os Itinerários \textbf{Formativos} ( \textbf{IF} ), com 1.200 horas, no mínimo.
Conforme preconizado nos documentos normativos nacionais, a FGB é composta por quatro áreas do conhecimento : Linguagens e suas Tecnologias, Matemática e suas Tecnologias, Ciências Humanas e Sociais Aplicadas e Ciências da Natureza e suas Tecnologias.
Os \textbf{IF}, além das quatro áreas citadas, abrangem a Formação Profissional e Técnica, sendo orientados para o aprofundamento e a ampliação das aprendizagens.
No Tocantins, as áreas do conhecimento da FGB serão efetivadas no cotidiano escolar por meio dos componentes curriculares já presentes no Currículo, enquanto os \textbf{IF} terão arranjos diversificados compostos por unidades curriculares.
Sendo assim, este Caderno apresenta possibilidades de Itinerários \textbf{Formativos} que devem ser desenvolvidas a partir de estratégias pedagógicas diversificadas, dentre elas, projetos, clubes, oficinas, núcleos de estudo, entre outras formas de trabalho.
Os estudantes poderão escolher as \textbf{Trilhas} de Aprofundamento a partir do plano individual de curso, da reflexão sobre seu Projeto de Vida, do diagnóstico das suas necessidades pedagógicas e dos seus interesses individuais e coletivos, de modo a aprofundar, ampliar e/ou acompanhar as aprendizagens.
Para elaboração das \textbf{Trilhas} de Aprofundamento, foram realizadas duas “ escutas ” voltadas para os estudantes do Ensino Médio, sendo uma realizada pela Seduc e a outra, em parceria com o Instituto PORVIR.
As escutas tinham como objetivo colher informações acerca do \textbf{perfil} e interesses dos estudantes do Ensino Médio.
Os resultados subsidiaram a escrita dos Itinerários \textbf{Formativos}.
Diante disso, apresentamos 17 ( dezessete ) \textbf{Trilhas} de Aprofundamento, a saber : - Aprofundar as aprendizagens relacionadas às competências gerais, às Áreas de Conhecimento e/ou à Formação Técnica e Profissional. - Consolidar a formação integral dos estudantes, desenvolvendo a autonomia necessária para que realizem seus projetos de vida. - Promover a incorporação de valores universais, como ética, liberdade, democracia, justiça social, pluralidade, solidariedade e \textbf{sustentabilidade}. - Desenvolver habilidades que permitam aos estudantes ter uma visão de mundo ampla e heterogênea, tomar decisões e agir nas mais diversas situações, seja na escola, seja no trabalho, seja na vida. \textbf{Trilhas} de Aprofundamento de Linguagens e suas Tecnologias - Amplifica!
A linguagem em movimento ; - Clube dos Literatos Juvenis ; - Eu sou o meu padrão! ; - Cultura Digital - na vibe das redes ; - Aperta o Play! \textbf{Trilhas} de Aprofundamento de Matemática e suas Tecnologias - Contribuições da matemática para o mundo \textbf{digital} ; - Como a Matemática se conecta com a Juventude, com a democracia e a sociedade? ; - Finanças Pessoais : o que o mundo exige na vida adulta que a gente pode aprender na escola? ; - Meu mundo, Meu futuro : Me ajuda a construir? - Modelagem Matemática aplicada à vida : construindo o saber matemático a partir das relações sociais. \textbf{Trilhas} de Aprofundamento de Ciências Humanas e Sociais Aplicadas - Vozes da juventude : passado e presente para um novo futuro ; - Sementes do cerrado : Cidadania e \textbf{Sustentabilidade} ; - Uma ideia na cabeça e uma câmera na mão. \textbf{Trilhas} de Aprofundamento de Ciências da Natureza e suas Tecnologias - Agronegócio e Agricultura Familiar ; - Ecoturismo em face do \textbf{empreendedorismo} ; - Energias Renováveis : Expectativa – Energia Fotovoltaica ( Solar ), Realidade – Usinas Elétricas ; - Nutrição e qualidade de vida : cuidando do corpo e da mente.
Ao escolher as \textbf{Trilhas} de Aprofundamento, o estudante poderá explorar as temáticas propostas, as quais trazem os quatro eixos \textbf{estruturantes} ( investigação científica, processos criativos, \textbf{mediação} e intervenção cultural, \textbf{empreendedorismo} ), os temas contemporâneos transversais, as habilidades gerais e específicas para a organização dos \textbf{IF} e a promoção do desenvolvimento humano na perspectiva dos valores universais.
O trabalho proposto é fruto da participação professores e técnicos pedagógicos da Seduc, contemplando suas práticas pedagógicas de profissionais da educação da rede pública de ensino do Tocantins e de agentes externos, como a Universidade Federal do Tocantins-UFT e Instituto Federal do Tocantins-IFTO.
Trata -se de material de apoio para as unidades escolares, a partir de objetivos e metodologias que reflitam a realidade local, o Projeto Político Pedagógico, as necessidades e características da comunidade escolar em consonância com o percurso do estudante, e a interação com seu Projeto de Vida.

% matched lemmas: digital, empreendedorismo, estruturante, formativo, if, mediação, perfil, sustentabilidade, trilha
\end{textsample}
